\section{Real campaigns, relays and randomised interventions}
\label{sec:world-full}

\paragraph{Documented campaigns.} The campaigns reported in 2026 occupy two corners of the frontier, neither of them the covert, self-written corner we studied in most detail. At one corner are high-volume campaigns. The largest Anthropic has reported ran more than 151 million exchanges between May and July 2026 through more than 3{,}500 accounts, peaking near 3 million a day---several hundred exchanges per account per day---and its accounts were linked because they all used one fixed prompt \citep{anthropic2026threat}. Campaigns like this are visible to per-account volume monitors and to template coordination, which our detectors caught almost every week (Section~\ref{sec:mechanism-full}). At the other corner are \emph{relays}: one forwarded almost 300{,}000 customer requests in ten days through 5{,}380 accounts, about six per account per day; another relayed more than 12 million exchanges in fourteen days \citep{anthropic2026threat}. Relays have low per-account volume and on-task content drawn from real users.

\paragraph{Why relays are the hard case.} Our covert campaigns were limited by supply: an attacker without a user base could not obtain or write enough prompts that are both useful and user-like. A relay has an unlimited supply of exactly such prompts. If a relaying account's traffic has the same joint law as the users it relays, Proposition~\ref{prop:law} leaves no query-only monitor able to separate it. That ``if'' is an assumption, not a measurement. A relay that interleaves many users' sessions within one account, or forwards only high-value requests, may differ in session coherence or in how topical each account is, and a monitor could in principle see that. We did not simulate relays; this is the most important experiment we leave open.

\paragraph{Beyond passive monitoring: the provider's own randomness.} Every threshold in this paper is learned from history, and history drifts and can be poisoned. A provider has one resource no attacker controls: its own randomness. If it throttles a flagged group in randomly chosen time blocks, it can test whether the work fails over to other accounts, with a null distribution that follows from the randomisation alone.

\begin{proposition}[Randomisation null]\label{prop:rand}
Fix, before any intervention, a flagged group $G$, a disjoint set $B$ of accounts to watch, and time blocks $1,\dots,L$. Draw a throttle schedule $Z$ for $G$ from a known design, independently of all traffic. For any statistic $V(Z,Y_B)$ of $B$'s traffic fixed in advance, the $p$-value $\Pr_{Z'}\{V(Z',Y_B)\ge V(Z,Y_B)\}$ over the design is valid under the null that $B$'s traffic does not depend on whether $G$ is throttled---whatever the benign traffic's law, drift or coordination.
\end{proposition}

This is Fisher's randomisation test \citep{fisher1935design,imbens2015causal} in a switchback design \citep{bojinov2023design}. In simulation on WildChat populations its null held on real, drifting traffic, and campaigns that reassigned throttled work to unflagged colluders were caught in every run, while those that did not lost 40\% of their throughput. But a benign team that shares a work queue and fails over legitimately fired the test 82--98\% of the time, throttling harms the legitimate accounts in a flagged group, and twelve weekly blocks take a quarter of a year. Whether such a test is specific in practice depends on how often benign groups respond to enforcement on other accounts---a rate only providers can measure. Appendix~\ref{app:randomized} gives details.
